% Generated draft: requires mathematical and source review.
\documentclass[11pt]{article}
\usepackage[margin=25mm]{geometry}
\usepackage{fontspec}
\setmainfont{Latin Modern Roman}
\usepackage{amsmath,amssymb,mathtools,unicode-math}
\setmathfont{Latin Modern Math}
\usepackage{enumitem,needspace,xurl,hyperref}
\setlength{\parindent}{0pt}
\setlength{\parskip}{5pt}
\setlength{\emergencystretch}{3em}
\newcommand{\N}{\mathbb{N}}
\newcommand{\Z}{\mathbb{Z}}
\newcommand{\Q}{\mathbb{Q}}
\newcommand{\R}{\mathbb{R}}
\newcommand{\C}{\mathbb{C}}
\newcommand{\F}{\mathbb{F}}
\newcommand{\A}{\mathbb{A}}
\newcommand{\PP}{\mathbb{P}}
\newcommand{\E}{\mathbb{E}}
\newcommand{\eps}{\varepsilon}
\newcommand{\dd}{\,\mathrm{d}}
\begin{document}
\section*{UnsolvedMath ID: 30006362}
\textbf{Source problem code:} catalogue entry
\subsection{Definitions and mathematical statement}
A factor is a von Neumann algebra with center consisting of scalar multiples of its identity, and a subfactor inclusion is a unital inclusion of factors. Let \(\Omega\) be a finite set. For an integer \(k\) with \(1\leq k\leq |\Omega|\), a permutation group \(G\leq\operatorname{Sym}(\Omega)\) is sharply \(k\)-transitive if, for any two ordered tuples \((x_1,\ldots,x_k)\) and \((y_1,\ldots,y_k)\), each having pairwise distinct entries in \(\Omega\), there is exactly one \(g\in G\) such that \(g x_i=y_i\) for every \(i\). Fix \(\omega\in\Omega\) and put
\[
H=G_\omega=\{g\in G:g\omega=\omega\}.
\]
An action \(\alpha:G\to\operatorname{Aut}(R)\) on a factor \(R\) is outer if \(\alpha_g\) is not an inner automorphism for every \(g\neq e\). Its restriction to \(H\) gives the crossed-product inclusion
\[
R\rtimes_{\alpha|_H}H\subset R\rtimes_\alpha G.
\]
In the usual finite-index subfactor setting, the principal graph is the rooted bipartite graph recording the multiplicities arising from the inclusion's iterated bimodule decompositions. The supplied record denotes the principal graph of a reference group--subgroup crossed-product inclusion by \(\Gamma_{G>H}\). It does not specify the reference factor, reference action, factor type, or precise principal-graph conventions; these data must be fixed in the intended setting before the assertion is fully determined.

The proposed assertion is the following. For a sharply \(k\)-transitive finite permutation group \(G\) and its point stabilizer \(H\), every subfactor inclusion \(Q\subset P\) whose principal graph is isomorphic to \(\Gamma_{G>H}\) is isomorphic, as an inclusion, to a crossed-product inclusion for this same pair \((G,H)\). Explicitly, there should exist a factor \(R\), an outer action \(\alpha:G\to\operatorname{Aut}(R)\), and a von Neumann algebra isomorphism
\[
\Phi:P\longrightarrow R\rtimes_\alpha G
\quad\text{such that}\quad
\Phi(Q)=R\rtimes_{\alpha|_H}H.
\]
The source does not specify the intended range of \(k\), beyond using sharply \(k\)-transitive groups, or the ambient class of subfactors quantified over. No factor type, irreducibility condition, or additional analytic hypothesis is imposed here to fill that gap. The hypothesis names a principal graph only; the supplied statement does not additionally require a dual principal graph or an identification of standard invariants.

The supplied excerpt asserts that the conjecture holds for \(k=2,3,4\). This is retained as a source claim, not as an independently verified theorem. Its background discussion also asserts that sharply \(3\)-transitive permutation groups form exactly two families, denoted \(L(q)\) and \(M(q)\), which are transitive extensions of groups denoted \(H(q)\) and \(S(q)\), respectively. It identifies \(L(q)=\operatorname{PGL}_2(q)\), describes its fractional-linear action on the projective line, and describes \(M(q)\) as having the same underlying set as \(\operatorname{PGL}_2(q)\), with an action whose branches depend on whether the determinant is a square or a nonsquare. The admissible parameters, the automorphism denoted \(\sigma\), and parts of the displayed formulas are missing or corrupted. The claimed example printed as \(S(32)\) and \(Z23\rtimes Q_8\) is likewise not unambiguously recoverable. These background assertions therefore do not supply additional precise hypotheses or constructions for the conjecture.
\subsection{Short English statement}
The conjecture asks whether the principal graph associated with a sharply \(k\)-transitive finite group and a point stabilizer forces every subfactor with that graph to arise from an outer action of the same group and subgroup. The supplied source claims the cases \(k=2,3,4\) and labels the general problem open. The intended range of \(k\), analytic setting, and some background notation require clarification.
\subsection{Sources}
\begin{itemize}
\item Local UnsolvedMath registry, supplied complete source record: permanent ID \(30006362\), exact source problem code \texttt{catalogue entry}, title ``Goldman-Type Theorem for Sharply Transitive Permutation Groups.'' The record belongs to the Oberwolfach Reports Open Problems collection and gives proposed year \(2025\); no individual proposer is supplied.
\item Supplied report citation: \emph{Subfactors and Applications (2026)}, Oberwolfach Reports, DOI \texttt{10.4171/owr/2025/34}, pp.\ 1823--1886. Supplied source URL: \url{https://doi.org/10.4171/owr/2025/34}. The title year \(2026\) and report designation \(34/2025\) are retained as supplied.
\item Supplied literature source: ``Oberwolfach Report 34/2025,'' \url{https://ems.press/content/serial-article-files/52239}.
\item Literature assessment dated exactly \texttt{2026-08-22}: the record labels the problem open and states that only the cases \(k=2,3,4\) are established in the record, leaving the general statement unresolved. This is supplied dated curation, not independently verified literature assessment. The supplied background expressly requests independent verification of time-sensitive status and cited claims.
\item The source record labels its statement \texttt{corrected\_verified} and reports that the official report was checked, the question isolated from surrounding exposition, and notation repaired. These are supplied verification claims; this draft does not claim to have visited or checked the report.
\item The original excerpt cites ``[4, Chapter XI, Section 2]'' for its background on sharply \(3\)-transitive groups. The bibliographic identity of reference [4] is not supplied.
\end{itemize}
\end{document}
